\clearpage
\Article{362}
With these principles established, we can easily determine the prime divisors of the function $x^\nu-1$. We assume that all those divisors which also divide some function
\SourceLeaf{241}
$x^\nu-1$, where $\nu'<\nu$, have already been found, and propose to investigate the remaining ones. They can all be included in the expression $(P,\rho^k)$, if $P$ is one of them and $k$ takes all values less than $\nu$ and coprime to it.

In Chapter \textsc{vi} we showed how the primitive roots of the equation $x^\nu=1$ can be divided into classes so that, when all are expressed as powers of one of them, the same division into classes results whichever primitive root is taken as the base; we called such collections of roots \emph{periods}. It is clear that the functions $x,x^\alpha,x^\beta,x^\gamma$, and so on, where $\alpha,\beta,\gamma$, and so on, denote all numbers coprime to $\nu$, can be divided into periods in the same way, and each larger period again into smaller ones, until periods of the form $x^k,x^{kp},x^{kpp}\dotstwo x^{kp^{m-1}}$ are finally reached. Once this is done, the following is clear:

$1$. Since every period is composed of smallest periods of the form $x^k+x^{kp}+\text{etc.}$, division by any prime function of degree $m$ gives a remainder which is a number.

$2$. Since all terms of a period can always be put in the form $x^{\kappa\pdot a^\alpha b^\beta c^\gamma}\dotstwo$, where $\kappa.\ a,b,c\dotstwo$ are fixed numbers, while all values may be substituted for $\alpha,\beta,\gamma\dotstwo$, it is clear that the period becomes itself when $x$ is replaced by $x^k$ and $k$ has the form $a^\alpha b^\beta c^\gamma\dotstwo(\mathrm{mod.}\,\nu)$. From this it is easy to see that all the functions $P,(P\;\rho^k)$, and so on, with $k$ denoting a number of this kind, give the same remainder when the period is divided by them.

$3$. Thus, after such a remainder has been subtracted, the period can be divided by the product of all the functions $(P,\rho^k)$.
